% preamble/theorems.tex
% ============================================
% 定理环境体系 —— 绛紫 + 香槟金
% 使用 \newtcbtheorem，每个环境接受 {标题}{标签} 两个参数
% 标题栏格式：环境名 + 编号 + "：" + 自定义小标题
%   调用示例：
%     \begin{theorem}{存在唯一性定理}{thm:exist}
%       设 $f$ 在区域 $G\subset\R\times\R^n$ 上连续且关于 $x$ 满足
%       Lipschitz 条件，则初值问题存在唯一解。
%     \end{theorem}
%     \begin{definition}{词项的定义}{def:term}
%       ……
%     \end{definition}
% ============================================

% ---- 全局正文定理样式：原书为无框、无色栏的行内标题 ----
\tcbset{
  qttheo/.style={
    breakable, enhanced, frame hidden, detach title,
    arc=0pt, boxrule=0pt, leftrule=0pt,
    fonttitle=\bfseries\theoremFont,
    fontupper=\theoremFont,
    colback=white, colbacktitle=white, coltitle=black,
    left=0pt, right=0pt, top=0pt, bottom=0pt,
    before skip=8pt, after skip=4pt,
    before upper={\tcbtitle\quad},
  }
}

% ---- 彩条标题格式：环境名 + 编号 + (有自定义标题时) "：" + 自定义标题 ----
\newcommand{\qttheotitle}[3]{%
  #1\ #2\IfBlankT{#3}{：#3}%
}

% ============================================================
% 类别 I — LogicBlue (绛紫主色)
% 定理 / 推论 / 命题 / 原理
% 编号采用章内计数（number within=chapter），与原书"定理 1.x"约定一致
% ============================================================
\newtcbtheorem[number within=section]{theorem}{定理}{
  qttheo,
  colback=white, colframe=LogicBlue,
  separator sign={},
  title={\qttheotitle{\thmt@thmname}{\thmt@thmnumber}{\thmt@tmptitle}},
}{thm}

\newtcbtheorem[use counter from=theorem]{corollary}{推论}{
  qttheo,
  colback=white, colframe=LogicBlue,
  separator sign={},
  title={\qttheotitle{\thmt@thmname}{\thmt@thmnumber}{\thmt@tmptitle}},
}{cor}

\newtcbtheorem[use counter from=theorem]{proposition}{命题}{
  qttheo,
  colback=white, colframe=LogicBlue,
  separator sign={},
  title={\qttheotitle{\thmt@thmname}{\thmt@thmnumber}{\thmt@tmptitle}},
}{prp}

\newtcbtheorem[use counter from=theorem]{prp}{原理}{
  qttheo,
  colback=white, colframe=LogicBlue,
  separator sign={},
  title={\qttheotitle{\thmt@thmname}{\thmt@thmnumber}{\thmt@tmptitle}},
}{prn}

% ============================================================
% 类别 II — LogicDarkBlue (深绛紫)
% 公理 / 问题 / 引理 / 注记
% ============================================================
\newtcbtheorem[use counter from=theorem]{axiom}{公理}{
  qttheo,
  colback=white, colframe=LogicDarkBlue,
  separator sign={},
  title={\qttheotitle{\thmt@thmname}{\thmt@thmnumber}{\thmt@tmptitle}},
}{ax}

\newtcbtheorem[use counter from=theorem]{problem}{问题}{
  qttheo,
  colback=white, colframe=LogicDarkBlue,
  separator sign={},
  title={\qttheotitle{\thmt@thmname}{\thmt@thmnumber}{\thmt@tmptitle}},
}{pb}

\newtcbtheorem[use counter from=theorem]{lemma}{引理}{
  qttheo,
  colback=white, colframe=LogicDarkBlue,
  separator sign={},
  title={\qttheotitle{\thmt@thmname}{\thmt@thmnumber}{\thmt@tmptitle}},
}{lem}

\newtcbtheorem[use counter from=theorem]{remark}{注记}{
  qttheo,
  colback=white, colframe=LogicDarkBlue,
  separator sign={},
  title={\qttheotitle{\thmt@thmname}{\thmt@thmnumber}{\thmt@tmptitle}},
}{rem}

% ============================================================
% 类别 III — GoldMain (香槟金)
% 定义 —— 金色边框 + 淡紫底 + 深紫标题文字
% ============================================================
\newtcbtheorem[use counter from=theorem]{definition}{定义}{
  qttheo,
  colback=PurplePale, colframe=GoldMain,
  fonttitle=\bfseries\theoremFont\color{PurpleDeep},
  separator sign={},
  title={\qttheotitle{\thmt@thmname}{\thmt@thmnumber}{\thmt@tmptitle}},
}{def}

% ============================================================
% 类别 IV — LogicLightBlue (浅绛紫)
% 构造 / 记号 / 说明
% ============================================================
\newtcbtheorem[use counter from=theorem]{construction}{构造}{
  qttheo,
  colback=white, colframe=LogicLightBlue,
  separator sign={},
  title={\qttheotitle{\thmt@thmname}{\thmt@thmnumber}{\thmt@tmptitle}},
}{con}

\newtcbtheorem[use counter from=theorem]{note}{记号}{
  qttheo,
  colback=white, colframe=LogicLightBlue,
  separator sign={},
  title={\qttheotitle{\thmt@thmname}{\thmt@thmnumber}{\thmt@tmptitle}},
}{no}

\newtcbtheorem[use counter from=theorem]{ill}{说明}{
  qttheo,
  colback=white, colframe=LogicLightBlue,
  separator sign={},
  title={\qttheotitle{\thmt@thmname}{\thmt@thmnumber}{\thmt@tmptitle}},
}{il}

% ============================================================
% 类别 V
% 示例 —— 无框正文，按节独立连续编号
% 调用方式仍为 \begin{example}{标题}{标签} ... \end{example}
% ============================================================
\newcounter{example}[section]
\renewcommand{\theexample}{\thesection.\arabic{example}}

\NewDocumentEnvironment{example}{m m}{%
  \par\medskip
  \refstepcounter{example}%
  \ifstrempty{#2}{}{\label{ex:#2}}%
  \noindent{\bfseries\theoremFont
    示例~\theexample\ifstrempty{#1}{}{：#1}}\quad
  \theoremFont
}{%
  \par\medskip
}

% ============================================================
%  解 (Solution) 环境（无编号）
% ============================================================
\newcounter{solution}
\renewcommand{\thesolution}{\arabic{solution}}
\NewDocumentEnvironment{solution}{}{%
  \par\proofFont
  \refstepcounter{solution}%
  \noindent\textbf{解.}\quad
}{%
  \hfill $\square$\par
}

% ============================================================
%  证明 (Proof) 环境（可自定义小标题，如 证明（充分性））
%  amsthm 已定义 proof，此处用标准语法重定义
% ============================================================
\renewenvironment{proof}[1][证明]{%
  \par\proofFont
  \noindent\textbf{#1}\quad
}{%
  \par
}

% ---- 行距 ----
\setstretch{1.2}

% ============================================================
%  附注 / 命名定理 / 命名引理（无编号，行内加粗，仿原书）
%  用于 附注 1/2/3、Cauchy-Peano 定理、Gronwall 引理 等
% ============================================================
\newcommand{\fuzhu}[2]{\par\medskip\noindent\textbf{附注~#1}\quad #2\par\medskip}
\newcommand{\namedthm}[2]{\par\medskip\noindent\textbf{#1~定理}\quad #2\par\medskip}
\newcommand{\namedlem}[2]{\par\medskip\noindent\textbf{#1~引理}\quad #2\par\medskip}
\newcommand{\namedcor}[2]{\par\medskip\noindent\textbf{推论~#1}\quad #2\par\medskip}
